\chapter*{Annexes}
\addcontentsline{toc}{chapter}{Annexes}

\section{Annexe A — Diagrammes UML}

Diagrammes de cas d'utilisation, de séquence, d'activité et de classes réalisés avec Draw.io / Mermaid. Voir figures~\ref{fig:use_case_diagram},~\ref{fig:sequence_diagram},~\ref{fig:activity_diagram} et~\ref{fig:class_diagram} (chapitres 2 et 3).

\section{Annexe B — Modèle de base de données}

Le schéma PostgreSQL est décrit par les migrations versionnées du dossier \texttt{backend/migrations/}, appliquées par des scripts dédiés. Le dictionnaire de données est présenté au tableau~\ref{tab:data_dictionary} (chapitre 3).

\section{Annexe C — Documentation des API}

Les points d'entrée REST sont regroupés par domaine. Extrait des principaux endpoints :

\begin{itemize}
    \item \texttt{POST /api/auth/login} — authentification (JWT) ;
    \item \texttt{POST /api/visits/register} — enregistrement patient + visite ;
    \item \texttt{POST /api/visits/:id/call} — appel du patient ;
    \item \texttt{POST /api/visits/:id/place} — affectation à une ressource ;
    \item \texttt{POST /api/visits/:id/start} — démarrage de la consultation ;
    \item \texttt{PATCH /api/visits/:id/priority} — triage / priorité ;
    \item \texttt{POST /api/visits/:id/close} — clôture avec outcome ;
    \item \texttt{GET /api/queue/worklists} — listes de travail ;
    \item \texttt{GET /api/resources/board} — état des ressources ;
    \item \texttt{GET /api/resources/eligible} — ressources affectables ;
    \item \texttt{POST /api/care/visits/:id/triage} — triage structuré ;
    \item \texttt{POST /api/care/visits/:id/transfer} — demande de transfert ;
    \item \texttt{GET /api/chief/*} — supervision (pipeline, SLA, alertes) ;
    \item \texttt{GET /api/events/stream} — flux temps réel ;
    \item \texttt{GET /api/ai/recommendations} — recommandations IA.
\end{itemize}

\section{Annexe D — Captures d'écran de la plateforme}

Les captures d'écran sont à déposer dans le dossier \texttt{rapport-stage/images/} avec les noms ci-dessous. Elles apparaissent automatiquement dans les figures du rapport (voir figure~\ref{fig:login_screen},~\ref{fig:dashboard_receptionniste},~\ref{fig:dashboard_infirmier},~\ref{fig:dashboard_medecin},~\ref{fig:dashboard_chef} et~\ref{fig:dashboard_admin}).

\begin{itemize}
    \item \texttt{login\_screen.png} — page de connexion ;
    \item \texttt{dashboard\_receptionniste.png} — dashboard Accueil ;
    \item \texttt{dashboard\_infirmier.png} — dashboard Infirmier (triage, plan des lits) ;
    \item \texttt{dashboard\_medecin.png} — dashboard Médecin ;
    \item \texttt{dashboard\_chef.png} — dashboard Chef de service (supervision) ;
    \item \texttt{dashboard\_admin.png} — dashboard Administrateur ;
    \item \texttt{architecture\_picu.png} — architecture en trois couches.
\end{itemize}

\section{Annexe E — Scénarios et commandes de test}

\paragraph{Invariants de ressources (\texttt{npm run test:v4})}
Vérifie qu'un lit en nettoyage ne peut être affecté, que la libération le bascule en nettoyage et que la transition directe vers « prêt » est refusée.

\paragraph{Règles de supervision (\texttt{npm run test:supervision})}
Vérifie qu'aucun seuil SLA non approuvé n'est activé.

\paragraph{Intelligence d'enregistrement (\texttt{npm run test:registration})}
Vérifie les indices d'enregistrement (motif à risque, ambulance, priorité élevée) sans modification de la priorité clinique.

\paragraph{Benchmark de triage}
\texttt{node benchmarks/run-triage-benchmark.js} mesure la latence du modèle NIM sur un jeu de cas et enregistre les résultats.

\section{Annexe F — Matrice de traçabilité complète}

\begin{table}[H]
\centering
\caption{Matrice Exigence $\rightarrow$ Implémentation $\rightarrow$ Test}
\label{tab:matrice_complete}
\small
\begin{tabular}{|p{3.8cm}|p{5cm}|p{4.4cm}|}
\hline
\textbf{Exigence} & \textbf{Implémentation} & \textbf{Validation} \\
\hline
Une visite active par patient & \texttt{visitService.registerVisitTransaction} & Conflit 409 \\
\hline
Un patient par ressource & \texttt{assertBedAssignable} + verrou & Test concurrence (T6) \\
\hline
Transitions autorisées & \texttt{visitService}, \texttt{careStage} & Tests par rôle \\
\hline
Ressource non affectable si nettoyage & \texttt{resourceService} & \texttt{test:v4} (T1) \\
\hline
Remise en service contrôlée & \texttt{setOperationalStatus} & \texttt{test:v4} (T3) \\
\hline
Seuils SLA approuvés uniquement & \texttt{supervisionRulesService} & \texttt{test:supervision} (T4) \\
\hline
Indices d'enregistrement & \texttt{registrationHints} & \texttt{test:registration} (T5) \\
\hline
Idempotence des actions & \texttt{findIdempotentActionResult} & Test idempotence (T7) \\
\hline
Approbation des actions IA & \texttt{action-policy} & Tests de politique IA \\
\hline
Traçabilité & \texttt{visit\_events} & Parcours E2E (T8) \\
\hline
\end{tabular}
\end{table}
